% !TEX root = ../main.tex
\chapter{Introduction}
\label{Chapter_Introduction}
Quantum dynamics in molecular and condensed-matter systems unfolds on
femtosecond timescales and is shaped by coherent couplings and
environmental fluctuations. In many molecular aggregates and related
excitonic systems, the relevant electronic degrees of freedom are embedded
in warm, structured surroundings that induce dephasing and relaxation
\cite{breuerpetruccione2007theoryopenquantum,
weiss2012quantumdissipativesystems}.
Understanding how these processes appear in measurable spectra is a central
problem in the spectroscopy of open quantum systems.

Linear absorption and emission identify transition energies and oscillator
strengths, but they compress several dynamical mechanisms into a
one-dimensional line shape \cite{mukamel1995principlesnonlinearoptical}.
Homogeneous dephasing, population relaxation, and static disorder can
therefore be difficult to separate from linear spectra alone.
Two-dimensional electronic spectroscopy (2DES) correlates excitation and
detection frequencies through sequences of ultrashort laser pulses
\cite{jonas2003twodimensionalfemtosecondspectroscopy,
gelzinisetal2019twodimensionalspectroscopynonspecialists}. It can separate
homogeneous and inhomogeneous broadening, resolve coupled transitions
through cross peaks, and track spectral evolution on femtosecond timescales
\cite{cho2009twodimensionalopticalspectroscopy}. On the theory side, the
measured signal can be expressed through multi-time response functions of
the dipole operator
\cite{mukamel1995principlesnonlinearoptical}.
In a broader physical context, 2DES is a key tool for studying
excitation-energy transport in systems containing several interacting
light-absorbing units, especially photosynthetic light-harvesting complexes
\cite{schlau-cohenetal2011twodimensionalelectronicspectroscopy,
lambrevetal2020insightsmechanismsdynamics}. In such systems, many
electronic transitions overlap spectrally, while the observed dynamics are
shaped by electronic coupling, environmental reorganisation, relaxation,
dephasing, and static disorder. This makes the additional frequency
resolution of 2DES especially useful for identifying transfer pathways,
excitonic couplings, and coherence signatures
\cite{schlau-cohenetal2011twodimensionalelectronicspectroscopy}. In
photosynthetic antenna complexes, 2DES has resolved
sub-\SI{100}{\femto\second} intraband dynamics together with
sub-picosecond to picosecond population transfer
\cite{zigmantasetal2006twodimensionalelectronicspectroscopy}. At the
same time, oscillatory signals require careful interpretation, because
long-lived beatings are not uniquely electronic and may instead have
vibrational or vibronic origin
\cite{lambrevetal2020insightsmechanismsdynamics}.

This thesis develops a computational workflow for simulating
two-dimensional electronic spectra of open quantum systems in the
weak-coupling regime. The aim is to connect reduced open-system dynamics
generated by master equations to time- and frequency-domain spectroscopic
observables.

\section{Scope and Approach}
\label{sec:intro_scope}

The reduced dynamics are described mainly by the Redfield master equation
\cite{breuerpetruccione2007theoryopenquantum,
weiss2012quantumdissipativesystems}.
This gives a microscopic weak-coupling link from a system--bath
Hamiltonian and bath correlation functions to relaxation and dephasing
terms. Because Redfield dynamics does not guarantee complete positivity,
the thesis is restricted to regimes where the weak-coupling and Markov
assumptions are reasonable
\cite{breuerpetruccione2007theoryopenquantum,
rivasetal2010markovianmasterequations}. For the monomer calibration
benchmark, a phenomenological Lindblad bath is also used.

On the spectroscopy side, the third-order response-function formalism
provides the conceptual framework
\cite{mancal2022nonlinearopticalspectroscopy}. Numerically, the signal is
constructed in the time domain by propagating the driven open-system
dynamics with the laser pulses included explicitly and by removing
lower-order contributions through discrete phase cycling
\cite{segarra-martietal2018accuratesimulationtwodimensional,
sunetal2025twodimensionalspectroscopyopen}. This avoids the impulsive-limit
approximation and retains finite-pulse and pulse-overlap effects relevant
to the benchmark comparisons.

\section{Model Systems and Physical Questions}
\label{sec:intro_models_questions}

The workflow is applied to a hierarchy of model systems that isolates
different spectroscopic signatures:
\begin{itemize}
\item \textbf{Monomer.} A single optical transition used to test the
	numerical workflow and to separate homogeneous broadening from static
	disorder.
\item \textbf{Dimer.} Two coupled sites that permit cross peaks, coherent
	excitonic mixing, and waiting-time dynamics.
\item \textbf{Larger aggregates.} Examples that test how the same workflow extends
	to denser spectra in more complex systems.
\end{itemize}

The central question is how relaxation, dephasing, static disorder, and
coherent coupling shape photon-echo behaviour, cross-peak structure, and
waiting-time dynamics in two-dimensional spectra.

\section{Thesis Outline}
\label{sec:intro_outline}

\autoref{chapter:oqs_open_quantum_systems} introduces the open-quantum-system
framework, derives the Redfield master equation, and discusses bath
correlation functions and spectral densities. \autoref{chapter:spctr_spectroscopy}
develops the nonlinear-spectroscopy formalism and shows how the reduced
open-system dynamics are converted into measurable optical observables.
\autoref{chapter:model_systems} specifies the Hamiltonians, dipole
operators, and bath assumptions for the model systems studied in this
thesis. \autoref{chapter:results} then presents the numerical results for
the monomer, the dimer, and the larger-system extension examples.

Technical supporting material is collected in the appendices. These cover the
Dyson-series derivation of the
\(n\)-th order polarisation (\autoref{app:response_functions}), the
Maxwell-side link from polarisation to signal
(\autoref{app:2dfts_maxwell}), ultrashort-pulse intuition for carrier-phase
effects (\autoref{app:phase_kick_intuition}), numerical
implementation details (\autoref{app:numerical_implementation}), and
supplementary comparison figures for the monomer and dimer benchmarks
(\autoref{app:supplementary_comparison_figures}).
